\chapter{Plans for Future Work}
\label{chap:further-works}

The results presented in the previous chapters indicate that the lateral
balancing problem is constrained not only by controller design, but also by the
physical actuation mechanism. In particular, the present moving mass produces
insufficient gravitational moment for rapid recovery, while the limited travel
of the LM2070 makes large transient rod motion unacceptable. Future work will
therefore proceed in two stages. The immediate plan is to increase the moving
mass and evaluate the resulting improvement experimentally. A longer-term plan
is to investigate a hybrid actuator that combines a moving mass with a reaction
wheel, allowing the two devices to address different time scales of the
balancing problem.


\section{Increasing the Moving Mass}
\label{sec:future-increasing-rod-mass}

The first planned modification is to increase the mass carried by the lateral
rod. The gravitational moment generated by the moving assembly is approximately

\begin{equation}
\tau_{g,\mathrm{rod}}
=
-m_{\mathrm{rod}}gr\cos\theta.
\label{eq:future-rod-gravity-torque}
\end{equation}

For the same rod displacement, increasing $m_{\mathrm{rod}}$ therefore
increases the available balancing moment. It also increases the relative
importance of the direct angular route, characterized near the upright
equilibrium by

\begin{equation}
\chi
=
\frac{m_{\mathrm{rod}}R^2}{J}.
\label{eq:future-direct-route-ratio}
\end{equation}

However, a larger moving mass also reduces the acceleration produced by a given
linear-motor force:

\begin{equation}
\left.\dot{u}_2\right|_F
=
\frac{F}{m_{\mathrm{rod}}}.
\label{eq:future-heavy-rod-acceleration}
\end{equation}

The expected improvement is therefore not monotonic. Additional mass provides
a stronger gravitational moment but also makes the rod slower and increases
the required actuator force. The purpose of the planned experiment is to
identify the useful range of moving mass under the actual force, travel, and
thermal limits of the LM2070.


\subsection{Modular Mass Assembly}

A modular mild-steel assembly has been designed so that mass can be added
symmetrically to the two ends of the rod. The design consists of two rod
holders, each with a nominal mass of $0.4~\mathrm{kg}$, and four removable mass
units, each with a nominal mass of $0.3~\mathrm{kg}$. The parts have an outside
diameter of $1.50~\mathrm{in}$ and use M6$\times$1.0 threaded connections. If
all six components are installed, the nominal added moving mass is

\begin{equation}
m_{\mathrm{add}}
=
2(0.4)+4(0.3)
=
2.0~\mathrm{kg},
\label{eq:future-total-added-mass}
\end{equation}

excluding fasteners and manufacturing variation. The removable design permits
the mass to be increased in several symmetric steps rather than in a single
large change.

\begin{figure}[htbp]
  \centering
  \includegraphics[width=0.92\textwidth]{06/rodHolder.png}
  \caption{CAD model of the mild-steel rod holder proposed for the
  increased-mass experiment.}
  \label{fig:future-rod-holder}
\end{figure}


The experimental configurations will be introduced progressively. A suitable
sequence is the original assembly, the two holders without additional units,
one mass unit on each side, and finally two mass units on each side. Symmetric
installation is necessary to avoid an unintended static offset and to reduce
asymmetric loading of the linear guide.

\section{Existing Lateral-Balancing Mechanisms}
\label{sec:future-existing-solutions}

The Wheelbot and Mini Wheelbot are representative modern platforms that use a
reaction wheel to control the unstable lateral degree of freedom
\cite{geist2022wheelbot,hose2025miniwheelbot}. The Wheelbot uses one wheel as
the ground-contact rolling wheel and an orthogonally mounted second wheel as a
reaction wheel as Figure~\ref{fig:wheelbot}. This arrangement directly applies a body torque and supports
rapid disturbance rejection and self-erection. The Mini Wheelbot follows the
same general reaction-wheel principle in a smaller platform and combines it
with learning-based and nonlinear model-predictive control.

\begin{figure}[htbp]
  \centering
  \includegraphics[width=0.42\textwidth]{06/wheelbot.png}
  \caption{Schematic of the Wheelbot's reaction wheel configuration.}
  \label{fig:wheelbot}
\end{figure}

The main alternatives relevant to the present project are summarized in
Table~\ref{tab:future-actuator-comparison}.

\begin{table}[htbp]
  \centering
  \caption{Comparison of lateral-balancing actuation mechanisms.}
  \label{tab:future-actuator-comparison}
  \small
  \begin{tabular}{p{0.13\textwidth}p{0.24\textwidth}p{0.23\textwidth}p{0.23\textwidth}}
    \hline
    Mechanism & Main torque source & Principal advantage & Principal limitation \\
    \hline
    Reaction wheel
    & $\tau_{\mathrm{rw}}=-I_f\dot{\omega}_f$
    & Fast and direct body torque without lateral travel
    & Finite motor torque, wheel speed, and angular-momentum capacity \\

    Moving mass
    & $\tau_{g,\mathrm{rod}}=-m_{\mathrm{rod}}gr\cos\theta$
    & Low-frequency gravitational moment without continued acceleration
    & Limited travel, slower response, and large transient displacement \\
    \hline
  \end{tabular}
\end{table}

A reaction wheel generates body torque by changing its angular momentum. If
the flywheel inertia is $I_f$ and its angular velocity is $\omega_f$, the
reaction torque applied to the remaining robot is

\begin{equation}
\tau_{\mathrm{rw}}
=
-I_f\dot{\omega}_f.
\label{eq:future-reaction-wheel-torque}
\end{equation}

This makes the reaction wheel particularly effective for fast transient
correction. However, a nonzero constant torque requires the wheel speed to
continue changing. Once the maximum wheel speed is reached, the wheel is
saturated and can no longer provide torque in the same direction without a
momentum-unloading action. The Wheelbot study also
notes that reaction-wheel balancing can require fast changes in motor direction
and high motor velocities \cite{geist2022wheelbot}.

The moving mass has the opposite characteristics. Its response is slower and
limited by the rod travel, but an offset mass can provide a persistent
gravitational moment without continued actuator acceleration. Strictly, this
is a sustained torque rather than continuous mechanical work: at a static
equilibrium the angular velocity is zero, so the ideal mechanical power is
zero. Nevertheless, the gravitational moment can remain present without the
unbounded speed increase required from a reaction wheel under constant torque.
The lateral-pendulum unicycle concept uses this gravitational coupling as its
primary stabilization mechanism \cite{daud2017lateralpendulum}.


\section{Hybrid Moving-Mass and Reaction-Wheel Actuator}
\label{sec:future-hybrid-actuator}

The complementary characteristics of the two actuators motivate a hybrid
design. A compact reaction wheel would be mounted on the moving rod assembly,
with its spin axis aligned to generate torque about the lateral lean axis. The
reaction-wheel motor and rotor would also contribute to the moving mass,
increasing the gravitational moment produced by a given rod displacement as shown in Figure~\ref{fig:hybrid-actuator}.

\begin{figure}[htbp]
  \centering
  \includegraphics[width=0.82\textwidth]{06/hybrid_actuator_concept.png}
  \caption{Hybrid actuator concept combining a moving mass with a reaction wheel.}
  \label{fig:hybrid-actuator}
\end{figure}

The reaction wheel would be used primarily during rapid transients, including
the beginning of a recovery maneuver and the reversal of the rod. The moving
mass would then provide the slower and sustained component of the balancing
moment. In this interpretation, the reaction wheel suppresses the instantaneous
rod overshoot, while the displaced mass supplies the longer-duration torque.

